% ===================== preambolo =====================
\documentclass[11pt,a4paper]{article}

\usepackage[T1]{fontenc}
\usepackage[utf8]{inputenc}
\usepackage[english]{babel}
\usepackage[a4paper,margin=25mm]{geometry}
\usepackage{amsmath,amssymb}
\usepackage{graphicx}
\usepackage{booktabs}
\usepackage{multirow}
\usepackage{xcolor}
\usepackage{caption}
\usepackage{enumitem}
\usepackage{siunitx}
\sisetup{per-mode=symbol,range-phrase={--},range-units=single}
\usepackage[hidelinks]{hyperref}

\captionsetup{font=small,labelfont=bf}
\setlist{itemsep=2pt,parsep=0pt,topsep=4pt}

% --- segnaposto per figure: compila senza immagini e si vede cosa manca ---
% Uso:  \figph{55mm}{fig:etichetta}{Didascalia}{che cosa deve mostrare}
\newcommand{\figph}[4]{%
  \begin{figure}[htbp]
    \centering
    \fcolorbox{gray!60}{gray!5}{%
      \begin{minipage}[c][#1][c]{0.92\linewidth}
        \centering\ttfamily\footnotesize\color{gray!80!black}
        [FIGURE PLACEHOLDER]\\[2pt]
        #4
      \end{minipage}}
    \caption{#3}
    \label{#2}
  \end{figure}}

% abbreviazioni dei controllori
\newcommand{\Co}{\textsf{C1}}
\newcommand{\Ct}{\textsf{C2}}
\newcommand{\Cth}{\textsf{C3}}

% =====================================================================
\title{\textbf{Terrain-Interaction Control of a Hexapod Robot:\\
a Measured Comparison of Three Controllers\\ on a Shared Simulation Plant}}

% TODO: inserire i nomi veri e le matricole
\author{[Student 1] \and [Student 2]}
\date{Field and Service Robotics --- A.Y. 2025/2026\\
Universit\`a degli Studi di Napoli Federico II\\[4pt]
\today}

\begin{document}
\maketitle

\begin{abstract}
\noindent
We compare three controllers for the PhantomX AX Metal Hexapod Mark~II on a
single high-fidelity Simscape Multibody plant: an open-loop kinematic tripod
controller (\Co), the torque-based terrain-search feedback of Arrigoni et
al.~\cite{arrigoni2024} (\Ct), and an attitude loop added on top of it
(\Cth). Each controller adds exactly one mechanism to the previous one, so
that every adjacent comparison isolates a single cause.

The headline outcome of the project was not in the original plan. Roughly half
of the effort went into \emph{finding and correcting two defects of our own
test bench}, both capable of inverting the conclusions and neither of which
announced itself: the robot walked, the animation was plausible and the
metrics came out. We report them in full, together with the measurement
discipline we adopted afterwards --- a measured noise floor, a readability
threshold fixed before looking at the data, and written retractions --- which
we consider the most transferable contribution of this work.

With that discipline in place, the measured result is that terrain search
\emph{costs} body attitude where there is nothing to search for and
\emph{repays} it on discontinuities, while the attitude loop gives that
attitude back and pays for it in energy. Neither mechanism changes the
response to a lateral impulse, and none of the three completes the
seven-obstacle course.
\end{abstract}

\tableofcontents
\newpage

% =====================================================================
\section{Introduction}
\label{sec:intro}

\subsection{The platform and the reference work}

The PhantomX AX Metal Hexapod Mark~II is an inexpensive, open-hardware
six-legged robot with eighteen Dynamixel AX-12A smart servos, three per leg
(coxa, femur, tibia). In its standard configuration it has \emph{no} force
sensors at the feet and no IMU: the only feedback available from the hardware
is the position and load estimate reported by each servo.

Arrigoni et al.~\cite{arrigoni2024} exploit precisely that constraint. Their
contribution is a feedback architecture that detects foot--ground contact
\emph{from motor torque alone}, by comparing the measured joint torque against
the torque predicted by an inverse-dynamics model of the robot, and uses the
detection to adapt the commanded foot height to terrain the trajectory
generator does not know about. The appeal of the idea is that it buys terrain
adaptation without adding a single sensor.

\subsection{Why we selected this project}

Three reasons made this an attractive final project rather than a
reimplementation exercise.

\begin{description}[style=nextline,leftmargin=0pt]
\item[The claim is falsifiable, and the paper does not quantify its price.]
The reference work demonstrates that terrain search improves behaviour on
unmodelled terrain. It does not report what the mechanism costs where there is
no terrain to discover. A controlled comparison on identical ground, with a
declared readability criterion, can answer that --- and the answer turned out
to be a real and systematic cost (Section~\ref{sec:results}).

\item[The contact detector is a model-based observer in disguise.]
The flag that drives the whole architecture is
$\mathrm{OR}_j\left(|\tau_j^{\mathrm{meas}} - \tau_j^{\mathrm{exp}}| >
\sigma_j\right)$, where $\tau^{\mathrm{exp}}$ comes from a rigid-body-tree
inverse-dynamics model. Its quality therefore depends on the \emph{accuracy of
a dynamic model}, which connects directly to the modelling and
state-estimation material of the course. We found out the hard way how
directly (Section~\ref{sec:estimator}).

\item[A hexapod on a shared plant makes the ablation clean.]
Because all three controllers drive the same multibody model, the same ground,
the same contact law and the same gait generator, each adjacent comparison
differs by one mechanism only. That is rarely available when comparing
published controllers against each other.
\end{description}

\subsection{What this report claims}

\begin{enumerate}
\item Two defects of the simulation bench, found, measured, corrected and
      documented, each of which had already produced published-looking numbers
      that were wrong (Section~\ref{sec:defects}).
\item A measurement protocol that makes a difference between two controllers
      \emph{declarable or not declarable}, rather than merely observed
      (Section~\ref{sec:method}).
\item A three-way comparison over seven tasks with that protocol applied
      (Section~\ref{sec:results}).
\item A feasibility check of the commanded torques against the real servo
      datasheet (Section~\ref{sec:feasibility}).
\end{enumerate}

% =====================================================================
\section{Experimental setup}
\label{sec:setup}

\subsection{The plant}

The robot is imported from the public URDF of the PhantomX description package
into Simscape Multibody: 25 rigid bodies (trunk plus 24 leg links), 18 revolute
joints, and six spherical contact elements at the feet. Joints are commanded in
\emph{position}; the resulting joint torque is therefore a reaction, not an
input --- a point that matters for Section~\ref{sec:feasibility}. Foot--ground
interaction uses a penalty contact law with Coulomb friction, tuned so that the
static penetration ($4.3$~mm in tripod stance) sits well above the transition
region ($2.0$~mm), where the force law would otherwise be ill-defined.

\figph{60mm}{fig:plant}{The Simscape Multibody plant: body, six three-degree-of-freedom legs, spherical contact elements and the terrain props used by the task battery.}{screenshot del modello Simscape / schema a blocchi dell'impianto}

\subsection{The three controllers}

\begin{table}[htbp]
\centering
\caption{The three controllers. Each adds one mechanism to the previous one.}
\label{tab:controllers}
\small
\begin{tabular}{@{}llp{85mm}@{}}
\toprule
 & \textbf{Plant} & \textbf{Mechanism} \\
\midrule
\Co  & \texttt{sim\_zero} & Open-loop kinematic tripod. Foot trajectories
  generated from the gait parameters, inverse kinematics per leg, joints
  commanded in position. No feedback of any kind. \\[3pt]
\Ct  & \texttt{sim\_zero} & \Co{} plus \emph{terrain search}: while a leg is in
  stance and the contact flag is low, the commanded foot height keeps
  descending until the torque residual exceeds the threshold. \\[3pt]
\Cth & \texttt{sim\_attitude} & \Ct{} plus an \emph{attitude loop}: two
  discrete PI regulators on body pitch and roll, whose outputs are distributed
  over the six feet as an additive height offset. \\
\bottomrule
\end{tabular}
\end{table}

The attitude loop acts \emph{in series, after} the terrain search: it reads the
already-corrected foot height and adds a term. We verified this by reading the
block connectivity inside the model file, not from documentation:

\begin{center}
\texttt{terrain\_search} $\rightarrow z_i \rightarrow \Sigma(+\Delta z_i)
\rightarrow$ \texttt{Saturation} $\rightarrow$ \texttt{Rate Limiter}
$\rightarrow$ leg
\end{center}

with
\begin{equation}
\Delta z = \mathbf{x}\,u_\theta + \mathbf{y}\,u_\phi ,
\label{eq:dz}
\end{equation}
where $u_\theta$ and $u_\phi$ are the pitch and roll regulator outputs and
$\mathbf{x},\mathbf{y}\in\mathbb{R}^6$ are the lever arms of the six feet.
Pitch uses $P=10$, $I=1$; roll uses $P=5$, $I=0.5$; both are discrete at
\SI{1}{\milli\second} with zero reference. Because \Cth{} contains \Ct{}, the
comparison $\Ct\rightarrow\Cth$ isolates the attitude loop, and
$\Co\rightarrow\Ct$ isolates terrain search.

\subsection{Gait parameters and where they come from}
\label{sec:gait}

The reference work delegates gait generation to NUKE and states explicitly
that it will not discuss it. The gait parameters therefore had to be chosen.
They are not arbitrary: three of five are fixed by a measured constraint, and
the remaining degree of freedom is a single one, because
\begin{equation}
T = \frac{S}{\beta_{\mathrm{stance}}\, v_{\mathrm{nom}}} .
\label{eq:gait}
\end{equation}

\begin{table}[htbp]
\centering
\caption{Gait parameters and their provenance.}
\label{tab:gait}
\small
\begin{tabular}{@{}llp{78mm}@{}}
\toprule
\textbf{Parameter} & \textbf{Value} & \textbf{Origin} \\
\midrule
phase vector & --- &
  Definition of the alternating tripod: groups
  $\{\mathrm{FL,MR,RL}\}$ and $\{\mathrm{FR,ML,RR}\}$. Not a tunable. \\
$\beta_{\mathrm{stance}}$ & $0.50$ &
  The only duty factor giving continuous support on \emph{exactly} three legs.
  Below it, instants with fewer than three feet exist (outside static
  stability); above it, six-foot windows appear in which 18 position
  constraints act on 6 degrees of freedom. \\
$S$ & \SI{0.06}{\meter} &
  Bounded by leg reach. At mid-stride the leg uses $80.9\%$ of
  $\ell_f+\ell_t=\SI{0.219}{\meter}$; the internal assertion trips at $85\%$,
  i.e.\ at $S=\SI{0.095}{\meter}$. The $85\%$ bound is itself measured: at
  $89\%$ the robot walked backwards. \\
$H$ & \SI{0.05}{\meter} &
  $1.5\times$ the measured obstacle height (\SI{32.8}{\milli\meter}, read from
  the stance-foot elevation in the single-obstacle run, not from the mesh). \\
$T$ & \SI{1.00}{\second} &
  Dependent variable, via~\eqref{eq:gait}. Gives
  $v_{\mathrm{nom}}=\SI{0.12}{\meter\per\second}$, verified at $83\%$ of the
  measured stability envelope (Section~\ref{sec:envelope}). \\
\bottomrule
\end{tabular}
\end{table}

The resulting Froude number is
$v^2/(g h) = 0.12^2/(9.81\times 0.154) = 0.0095$: inertial forces are one
percent of gravity. The march is firmly quasi-static, which is what legitimises
the \emph{static} stability assumption the tripod gait is built on.

\subsection{Task battery}

\begin{table}[htbp]
\centering
\caption{The seven tasks.}
\label{tab:tasks}
\small
\begin{tabular}{@{}llp{80mm}@{}}
\toprule
\textbf{ID} & \textbf{Name} & \textbf{What it exercises} \\
\midrule
T2  & flat, speed sweep & Baseline, and the stability envelope in speed. \\
T3  & curved path       & Yaw command \SI{0.1}{\radian\per\second}. \\
T4  & \SI{8}{\degree} ramp & Continuous slope: body attitude versus terrain. \\
T4D & \SI{8}{\degree}/\SI{8}{\degree} bump & Up-then-down slope, two angular interfaces. \\
T5  & single obstacle   & One unmodelled step, \SI{32.8}{\milli\meter}. \\
T6  & seven obstacles   & Repeated unmodelled discontinuities, half-width so
                          that left and right sides are raised alternately. \\
T7  & lateral impulse   & Momentum step from $1\%$ to $1600\%$ of nominal. \\
\bottomrule
\end{tabular}
\end{table}

\figph{65mm}{fig:tasks}{The terrain configurations of the task battery.}{montaggio dei terreni: piano, rampa, dosso, ostacolo singolo, percorso a sette ostacoli}

% =====================================================================
\section{Methodology: when is a difference a result?}
\label{sec:method}

Three rules, each adopted after having violated it at least once. We regard
them as the most transferable part of this work, because they are what
separated the results in Section~\ref{sec:results} from the numbers we
believed before.

\subsection{A measured noise floor, and a readability threshold}
\label{sec:noise}

Every run is repeated three times with the stance height $z_0$ displaced by
$\pm\SI{0.2}{\milli\meter}$ --- a quantity that has nothing to do with the
controller and everything to do with where the ground happens to be. The
resulting excursion of each metric is the \emph{noise floor} of that metric,
on that task, for that controller.

\begin{quote}
A difference between two controllers is declared only if it is at least
\textbf{three times} the worse of the two noise floors.
\end{quote}

This is not a formality. Applying it removed several differences that looked
like results; applying it \emph{badly} --- with a floor measured on a
different controller --- made us believe for a day that nothing at all was
declarable. The floor is measured on four tasks (flat, curve, bump, obstacle)
for all three controllers, each against itself.

Throughout Section~\ref{sec:results}, a figure in parentheses such as
$(16.3\times)$ is the ratio of the observed difference to that noise floor.
Entries marked \textsf{n.c.} (not conclusive) fall below $3\times$; entries
marked \textsf{n.m.} have no measured floor on that task and are reported
without a verdict.

\subsection{The criterion is written before the run}

Every diagnostic script in the repository carries, in its header, the question
it asks, the threshold it will apply, and the possible conclusions --- all
written before the data exist. This matters because the alternative is
available and tempting: on this project a readability threshold calibrated
after seeing the data would have moved the stability limit from $1.0\times$ to
$1.2\times$ depending on which multiple we chose.

\subsection{Retractions stay written}

Every claim we withdrew is marked \texttt{[RETRACTED \textless date\textgreater]}
in place, with the reason. We withdrew at least five during the project. On
this work a wrong claim survived three days precisely because it was
plausible, and the discipline exists to make that visible rather than tidy.

% =====================================================================
\section{Two defects of the test bench}
\label{sec:defects}

The largest single block of work in this project was not the comparison. It
was discovering that the bench we were measuring with was wrong in two places,
and that the second defect was \emph{caused by fixing the first}.

\subsection{The URDF inertias are wrong by a factor of about one thousand}
\label{sec:inertias}

The plant takes masses and inertias from a public URDF. The masses are
correct; the inertia tensors are not. Two independent arguments, neither of
which depends on trusting any particular source for the file:

\begin{description}[style=nextline,leftmargin=10pt]
\item[Radius of gyration.] From $I = m r^2$, the declared trunk inertia places
the trunk's mass at $r = \SI{1.79}{\meter}$ from its own centre. The trunk is
\SI{25}{\centi\meter} across. Every link gives $r = \SI{0.46}{\meter}$; the
links are \SIrange{2}{12}{\centi\meter} long. Read in
$\mathrm{kg}\cdot\mathrm{m}^2$, as the URDF standard requires, those numbers
describe an object that cannot exist.
\item[All 24 links carry identical inertia,] to the same decimal. Coxa, femur
and tibia differ in shape and length. It is a placeholder, not a measurement.
\end{description}

Dividing by $10^3$ brings everything into the physically plausible range,
which is the signature of a CAD export performed with masses in grams. The
effect on conclusions, not merely on values, is large: with the URDF inertias
the pitch of \Co{} on rough ground reached \SI{45}{\degree}, against
\SI{15}{\degree} with corrected ones, and the cost ratio between \Ct{} and
\Co{} moved from $1.35$ to $0.68$ --- it \emph{changed sign}.

\subsection{Correcting the robot broke the contact detector}
\label{sec:estimator}

The contact flag of \Ct{} is not a threshold on measured torque. The model
computes
\begin{equation}
\mathrm{contact} = \bigvee_{j\in\{c,f,t\}}
\left( \left| \tau_j^{\mathrm{meas}} - \tau_j^{\mathrm{exp}} \right| > \sigma_j \right),
\label{eq:flag}
\end{equation}
and $\tau^{\mathrm{exp}}$ is produced by an inverse-dynamics block carrying a
\texttt{rigidBodyTree} imported \emph{from the same URDF}.

We corrected the inertias of the simulated robot only. From that moment the
robot and the model predicting its torque were two different robots, with
inertias differing by three orders of magnitude. The residual
$|\tau^{\mathrm{meas}} - \tau^{\mathrm{exp}}|$ became large \emph{always},
with or without contact, and the flag stayed high over $95\%$ of the swing
phase. With the flag stuck at one, the terrain search never enters the branch
that grows the foot extension, and the command degenerates into a constant.

\begin{quote}
For three days, \Ct{} did not search the terrain. Nothing signalled it: the
robot walked, the animation was plausible, the metrics came out. Every \Ct{}
row produced in that window described a broken controller.
\end{quote}

\paragraph{How we measured it.} A single ten-second run, in which the log
contains simultaneously the measured torque, the expected torque and the
\emph{normal force at the foot}. The flag can then be recomputed offline for
any threshold and compared against the force, which is ground truth.

\begin{table}[htbp]
\centering
\caption{Contact flag against measured foot force, at unchanged threshold.
The second row is the diagnosis: a flag that is always high is wrong exactly
as often as ours was, so ours carried no information.}
\label{tab:estimator}
\small
\begin{tabular}{@{}lcc@{}}
\toprule
 & \textbf{estimator misaligned} & \textbf{aligned} \\
\midrule
flag error during stance                & $7.0\%$ & $\mathbf{0.6\%}$ \\
error of an \emph{always-high} flag     & $7.0\%$ & $3.2\%$ \\
does the flag carry information?        & \textbf{no} & yes \\
swings with no reset of the search      & $20/60$ & $2/60$ \\
\bottomrule
\end{tabular}
\end{table}

\paragraph{The fix.} A routine rewrites the mass and inertia of each body of
the estimator's rigid-body tree with the same corrected values applied to the
plant, forces the masked subsystem to re-evaluate, and verifies in three
successive ways --- write, fallback, read-back --- stopping with an error if
any of them fails. A silent failure here would restore exactly the problem it
closes. It is called immediately after the inertia correction in every task
script, in memory, never saving the model.

\subsection{Re-tuning the contact threshold}
\label{sec:threshold}

With the estimator aligned the flag works, but at the original threshold
$[0.04,\,0.1,\,0.1]$ it still stays high over $41.9\%$ of swing, and $39.7\%$
of that is the coxa alone: its $p_{90}$ in swing ($0.0421$) overlaps its
$p_{10}$ in stance ($0.0469$). This has a physical explanation --- the coxa
rotates about the vertical axis while contact is vertical, so contact
transmits almost no torque to it, while in swing it carries the whole lateral
acceleration of the leg.

The adopted triple is $[0.385,\,0.136,\,0.105]$, which sets the coxa threshold
above its own stance median, i.e.\ switches it off. \textbf{This is a change of
rule, not a re-tuning:} the flag becomes an OR over two joints. It was decided
on the obstacle course, where the outcome is binary and not subject to
measurement noise, and it must be declared as our choice: it is no longer the
value present in the reference implementation.

\subsection{A slew-rate limiter, and why it belongs in the plant}
\label{sec:slew}

The attitude model originally carried a rate limiter of
$\pm\SI{0.3}{\meter\per\second}$ on each foot height command, which the
open-loop plant did not have. That is an asymmetry: the limiter constrains
\emph{every} height command, including the one produced by terrain search, so
a \Ct{}--\Cth{} comparison would have confounded the attitude loop with its own
limiter.

We therefore placed the same limiter, at
$\pm\SI{0.4}{\meter\per\second}$, in both plants. The value is not free: the
nominal gait peaks at \SI{0.375}{\meter\per\second} of vertical foot speed, so
at $\pm0.4$ the nominal march is untouched and the limiter acts only on fast
commands --- terrain search and attitude corrections. We verified that the
\Co{} rows and the \Co{} noise floor are byte-identical with and without the
limiter on five of the seven tasks and on the nominal cell of the sixth.

This turned out to matter more than expected: the limiter removed the lateral
drift of \Ct{} on the obstacle course (from $0.296$ to \SI{0.104}{\meter}) and
an anomalous noise sensitivity of the energy of \Cth{}. An implementation
detail, not a controller, with effects larger than some of the differences
between controllers.

% =====================================================================
\section{Results}
\label{sec:results}

All rows below come from a single campaign run after the corrections of
Section~\ref{sec:defects}, with the slew-rate limiter present on all three
plants: seven tasks $\times$ three controllers, plus a noise floor of four
tasks $\times$ three runs $\times$ three controllers. Verdicts in bold pass
the $3\times$ rule of Section~\ref{sec:noise}.

\subsection{Flat ground and curved path: terrain search costs attitude}

\begin{table}[htbp]
\centering
\caption{Flat ground (T2) and curved path (T3), nominal speed. Ratio to the
noise floor in parentheses; \textsf{n.c.} = below $3\times$.}
\label{tab:t2t3}
\small
\begin{tabular}{@{}lrrrll@{}}
\toprule
 & \Co & \Ct & \Cth & $\Co\rightarrow\Ct$ & $\Ct\rightarrow\Cth$ \\
\midrule
\multicolumn{6}{@{}l}{\emph{T2 --- flat, nominal speed}}\\
pitch max [deg]      & 0.514 & 1.95  & 0.356 & \textbf{+279\%} (16.3) & \textbf{$-$82\%} (18.1) \\
roll max [deg]       & 0.206 & 0.280 & 0.205 & +36\% (n.c.)           & $-$27\% (n.c.) \\
cost of transport    & 0.888 & 0.865 & 1.11  & $-$3\% (n.c.)          & \textbf{+28\%} (27.7) \\
energy [J]           & 18.6  & 18.2  & 23.3  & $-$2\% (n.c.)          & \textbf{+28\%} (28.4) \\
torque RMS [N\,m]    & 0.400 & 0.403 & 0.372 & +1\% (n.c.)            & \textbf{$-$8\%} (8.0) \\
peak torque [N\,m]   & 1.87  & 2.02  & 2.23  & +8\% (n.c.)            & +11\% (n.c.) \\
mean speed [m/s]     & 0.135 & 0.136 & 0.135 & \textbf{+1\%} (6.9)    & $-$0\% (n.c.) \\
\midrule
\multicolumn{6}{@{}l}{\emph{T3 --- curve, $+0.1$ rad/s}}\\
pitch max [deg]      & 0.579 & 1.85  & 0.392 & \textbf{+219\%} (40.0) & \textbf{$-$79\%} (45.9) \\
roll max [deg]       & 0.224 & 0.351 & 0.210 & \textbf{+57\%} (4.6)   & \textbf{$-$40\%} (5.1) \\
cost of transport    & 0.970 & 0.951 & 1.22  & \textbf{$-$2\%} (3.6)  & \textbf{+28\%} (24.3) \\
energy [J]           & 30.3  & 29.5  & 37.8  & \textbf{$-$3\%} (6.4)  & \textbf{+28\%} (22.5) \\
torque RMS [N\,m]    & 0.398 & 0.389 & 0.376 & \textbf{$-$2\%} (7.0)  & \textbf{$-$3\%} (7.3) \\
peak torque [N\,m]   & 1.91  & 2.11  & 2.43  & +10\% (n.c.)           & +15\% (n.c.) \\
\bottomrule
\end{tabular}
\end{table}

On ground that is exactly what the trajectory generator believes it to be,
terrain search makes the body attitude \emph{worse} by a factor of three to
four, and buys nothing: energy and cost of transport are unchanged within
noise on flat ground, and improve by two to three percent on the curve. The
mechanism keeps working where there is nothing to find, and the body pays for
it.

\figph{70mm}{fig:pitch-flat}{Body pitch over two gait cycles on flat ground, for the three controllers. The oscillation of \Ct{} is the terrain search acting on ground it already knows.}{grafico beccheggio vs tempo, T2, tre curve sovrapposte}

\subsection{Discontinuities: terrain search repays}

\begin{table}[htbp]
\centering
\caption{Bump (T4D) and single obstacle (T5).}
\label{tab:t4dt5}
\small
\begin{tabular}{@{}lrrrll@{}}
\toprule
 & \Co & \Ct & \Cth & $\Co\rightarrow\Ct$ & $\Ct\rightarrow\Cth$ \\
\midrule
\multicolumn{6}{@{}l}{\emph{T4D --- bump, \SI{8}{\degree}/\SI{8}{\degree}}}\\
pitch max [deg]    & 8.65  & 10.4  & 1.98  & \textbf{+21\%} (19.8)  & \textbf{$-$81\%} (93.8) \\
roll max [deg]     & 2.12  & 1.49  & 0.611 & \textbf{$-$30\%} (6.4) & \textbf{$-$59\%} (7.6) \\
cost of transport  & 1.12  & 1.04  & 1.45  & \textbf{$-$7\%} (5.1)  & \textbf{+40\%} (16.8) \\
energy [J]         & 64.1  & 59.9  & 83.6  & \textbf{$-$7\%} (3.5)  & \textbf{+40\%} (18.8) \\
mean speed [m/s]   & 0.130 & 0.133 & 0.131 & \textbf{+2\%} (3.6)    & \textbf{$-$1\%} (3.0) \\
\midrule
\multicolumn{6}{@{}l}{\emph{T5 --- single obstacle, \SI{32.8}{\milli\meter}}}\\
pitch max [deg]    & 7.56  & 6.41  & 1.73  & \textbf{$-$15\%} (16.4)& \textbf{$-$73\%} (19.8) \\
roll max [deg]     & 4.49  & 2.31  & 1.13  & $-$49\% (n.c.)         & $-$51\% (n.c.) \\
cost of transport  & 1.20  & 1.08  & 1.40  & \textbf{$-$10\%} (5.2) & \textbf{+30\%} (17.1) \\
energy [J]         & 48.1  & 44.3  & 57.5  & \textbf{$-$8\%} (3.6)  & \textbf{+30\%} (17.2) \\
torque RMS [N\,m]  & 0.408 & 0.403 & 0.380 & $-$1\% (n.c.)          & \textbf{$-$6\%} (9.9) \\
peak torque [N\,m] & 4.46  & 8.33  & 4.01  & +87\% (n.c.)           & $-$52\% (n.c.) \\
obstacle cleared   & yes   & yes   & yes   & same                   & same \\
\bottomrule
\end{tabular}
\end{table}

Here the picture reverses, and this is the result that reproduces the claim of
the reference work: on a discontinuity the search improves roll, energy, cost
of transport and task completion. One solid exception remains on the bump:
\textbf{peak pitch is $21\%$ worse}, at a ratio of $19.8$. On a continuous
slope the search extends the downhill feet and the body follows the gradient
--- it adapts to the terrain rather than resisting it.

\figph{70mm}{fig:obstacle}{Body attitude during the obstacle crossing (T5). Shaded: the window in which at least one foot is on the obstacle, identified kinematically.}{grafico assetto vs tempo su T5, con finestra del passaggio evidenziata}

\subsection{The ramp: what the attitude loop is for}

\begin{table}[htbp]
\centering
\caption{\SI{8}{\degree} ramp (T4). Noise floor not measured on this task;
no verdicts are issued, but the behaviour is qualitative and unambiguous.}
\label{tab:t4}
\small
\begin{tabular}{@{}lrrr@{}}
\toprule
 & \Co & \Ct & \Cth \\
\midrule
body attitude minus ramp angle [deg] & 0.106 & 0.259 & $-7.54$ \\
peak-to-peak pitch on the ramp [deg] & 0.994 & 1.97  & 0.839 \\
roll on the ramp [deg]               & 0.662 & 0.563 & 0.314 \\
cost of transport                    & 1.15  & 1.09  & 1.58 \\
climb completed                      & yes   & yes   & yes \\
\bottomrule
\end{tabular}
\end{table}

\Co{} and \Ct{} align the body with the ramp, to within a few tenths of a
degree. \Cth{} holds it $7.54^\circ$ away from the ramp, i.e.\ within
$0.5^\circ$ of \emph{horizontal}, which is exactly what a loop with zero
attitude reference is built to do --- and it is also the behaviour the
reference work states as its own objective. It costs $45\%$ in cost of
transport.

\figph{70mm}{fig:ramp}{Body pitch while climbing the \SI{8}{\degree} ramp. \Co{} and \Ct{} follow the slope; \Cth{} holds the trunk horizontal.}{grafico beccheggio vs tempo su T4, tre curve, con la pendenza della rampa come riferimento}

\subsection{Obstacle course: the comparison does not discriminate}

\begin{table}[htbp]
\centering
\caption{Seven-obstacle course (T6). Only the binary outcome and its validity
are reported.}
\label{tab:t6}
\small
\begin{tabular}{@{}lrrr@{}}
\toprule
 & \Co & \Ct & \Cth \\
\midrule
distance travelled [m]      & 2.91  & 3.12  & 2.99 \\
max lateral drift [m]       & 0.067 & 0.104 & 0.058 \\
course completed, valid     & no    & no    & no \\
\bottomrule
\end{tabular}
\end{table}

``Completed'' counts only if the maximum lateral drift stays below
\SI{0.183}{\meter}; beyond that the robot may have missed an obstacle entirely
rather than stepped over it. The threshold comes from the obstacle geometry,
computed independently of any controller.

\textbf{None of the three completes the course.} The fine metrics of this task
must not be quoted: the three controllers follow different trajectories and
therefore do not meet the same obstacles. On this task the comparison does not
separate the controllers, and saying so is the result.

\subsection{Lateral impulse: a gap common to all three}

\begin{table}[htbp]
\centering
\caption{Residual lateral deviation after an impulsive disturbance (T7).}
\label{tab:t7}
\small
\begin{tabular}{@{}lrrr@{}}
\toprule
 & \Co & \Ct & \Cth \\
\midrule
residual deviation, $1\times$ impulse [mm]  & 1.34 & 6.04 & 0.164 \\
residual deviation, $16\times$ impulse [mm] & 393  & 400  & 392 \\
recovers at $16\times$                      & no   & no   & no \\
\bottomrule
\end{tabular}
\end{table}

At large impulse the three curves coincide to within $2\%$: none of the three
recovers a measurable fraction of the deviation. Contact detection is not the
mechanism in play --- this is an architectural gap shared by all three, and the
clearest motivation for any future controller on this platform.

\subsection{The stability envelope in speed}
\label{sec:envelope}

A sweep of six speed factors, judged on body bounce with a threshold declared
beforehand, locates the envelope of the open-loop gait.

\begin{table}[htbp]
\centering
\caption{Body bounce versus commanded speed factor, \Co{}. Threshold
\SI{10}{\milli\meter}, declared before the sweep.}
\label{tab:envelope}
\small
\begin{tabular}{@{}lrl@{}}
\toprule
\textbf{factor} & \textbf{bounce [mm]} & \\
\midrule
1.00 & 3.2  & nominal operating point \\
1.20 & 3.4  & last cell with stable body motion \\
1.30 & 9.4  & at the threshold \\
1.40 & 19.7 & fails \\
\bottomrule
\end{tabular}
\end{table}

\figph{65mm}{fig:envelope}{Speed sweep: body bounce and task fraction versus commanded speed factor.}{grafico rimbalzo del corpo e frazione del task vs fattore di velocita'}

% =====================================================================
\section{Feasibility on the real actuators}
\label{sec:feasibility}

The campaign runs with ideal actuators: the simulator grants any torque. The
AX-12A stalls at \SI{1.5}{\newton\meter}, and the manufacturer recommends
operating at one fifth of stall for stable motion. Since the joints are
commanded in position, torque is a reaction rather than an input, so what
follows measures what each controller \emph{demands}, not how it would behave
with real motors.

\begin{table}[htbp]
\centering
\caption{Commanded torque against the AX-12A limit. RMS is the worst-loaded
joint of each run, as a fraction of stall; ``peak'' is the maximum torque as a
multiple of stall; ``over'' is the fraction of samples above stall.}
\label{tab:feasibility}
\small
\begin{tabular}{@{}lrrrrrrrrr@{}}
\toprule
& \multicolumn{3}{c}{RMS / stall} & \multicolumn{3}{c}{peak / stall}
& \multicolumn{3}{c}{samples over} \\
\cmidrule(lr){2-4}\cmidrule(lr){5-7}\cmidrule(lr){8-10}
task & \Co & \Ct & \Cth & \Co & \Ct & \Cth & \Co & \Ct & \Cth \\
\midrule
T2  & 46\% & 44\% & 39\% & 1.2 & 1.3 & 1.5 & 0.40\% & 0.45\% & 0.38\% \\
T3  & 46\% & 42\% & 39\% & 1.3 & 1.4 & 1.6 & 0.50\% & 0.40\% & 0.43\% \\
T4  & 44\% & 44\% & 45\% & 1.4 & 2.3 & 1.8 & 0.39\% & 0.49\% & 0.48\% \\
T4D & 42\% & 42\% & 39\% & 1.7 & 2.7 & 2.1 & 0.45\% & 0.41\% & 0.49\% \\
T5  & 45\% & 44\% & 40\% & 3.0 & 5.6 & 2.7 & 0.61\% & 0.50\% & 0.51\% \\
T6  & 42\% & 43\% & 42\% & 5.0 & 4.9 & 2.8 & 0.91\% & 0.67\% & 0.81\% \\
T7  & 46\% & 44\% & 39\% & 1.2 & 1.3 & 1.5 & 0.41\% & 0.45\% & 0.38\% \\
\bottomrule
\end{tabular}
\end{table}

Two regimes, and they must be read separately:

\begin{itemize}
\item \textbf{The march exceeds the recommended load.} The worst-loaded joint
sits between $39\%$ and $46\%$ of stall on every task and for every
controller --- $2.3\times$ the one fifth of stall the manufacturer indicates
for stable motion. It stays below stall, but not within the recommended
envelope.
\item \textbf{The impact peaks leave stall}, up to $5.0\times$ on the obstacle
course, but on at most $0.91\%$ of samples: these are isolated events, one per
foot strike, lasting tens of milliseconds.
\end{itemize}

\Cth{} consistently demands the lowest continuous torque ($39$--$45\%$) and the
lowest impact peaks on the tasks with discontinuities, which is consistent with
an attitude loop that absorbs part of the impact in body motion rather than in
the joints.

\figph{65mm}{fig:feasibility}{Commanded torque against the AX-12A limit: continuous (RMS of the worst joint) and peak, per task and controller.}{grafico a due pannelli: RMS/limite e picco/limite per task, tre controllori}


% =====================================================================
\section{Discussion}
\label{sec:discussion}

\subsection{The two mechanisms do different things, and only one of them pays}

Reading Tables~\ref{tab:t2t3}--\ref{tab:feasibility} along the two adjacent
comparisons gives a consistent picture across all seven tasks.

\paragraph{Terrain search alone (\Co$\rightarrow$\Ct) is not an attitude
mechanism.} On the two tasks where the ground is exactly what the trajectory
generator assumes --- flat (T2) and curved (T3) --- it makes pitch three to
four times worse, by a margin $16$ to $40$ times the noise floor, and returns
nothing measurable in energy. On the tasks with real terrain it does better:
roll improves by $30\%$ on the bump and $49\%$ on the obstacle, cost of
transport drops $7$--$10\%$, and the mean speed rises by $2$--$3\%$. But the
pitch, which is the degree of freedom the search acts on, still degrades by
$21\%$ on the bump.

The reason is visible in the mechanism rather than in the numbers. The search
lowers a swinging foot until the torque residual declares contact; the
declaration is made per leg, asynchronously, and the resulting height
correction is a step. On ground the generator already models correctly the
steps are spurious, and three legs stepping at slightly different instants is
exactly a pitch excitation. Arrigoni et al.~\cite{arrigoni2024} introduce the
mechanism for terrain the robot cannot see; our measurement says that on
terrain it \emph{can} see, the mechanism is a net cost, and that cost does not
disappear when the terrain becomes irregular --- it is simply joined by a
benefit on the other axes.

\paragraph{Attitude feedback (\Ct$\rightarrow$\Cth) is the mechanism that buys
attitude, and it buys it with energy.} The sign is the same on every task and
every magnitude clears the $3\times$ rule by a wide margin: pitch
$-82\%$ (flat), $-79\%$ (curve), $-81\%$ (bump), $-73\%$ (obstacle), with
noise ratios from $18$ to $94$. On the $8^\circ$ ramp it is qualitative rather
than incremental: \Co{} and \Ct{} carry the trunk parallel to the ramp
($0.11^\circ$ and $0.26^\circ$ of body-minus-ramp angle), while \Cth{} holds it
within half a degree of \emph{horizontal} ($-7.54^\circ$ of body-minus-ramp on
an $8^\circ$ slope). Those are two different behaviours, not two qualities of
the same behaviour, which is why that row carries no percentage verdict.

The price is equally consistent: $+28\%$ energy and cost of transport on flat
and curved ground, $+40\%$ on the bump and the ramp, $+30\%$ on the obstacle,
all at $17$ to $28$ times the noise floor. Interestingly the \emph{continuous}
torque goes down ($-3$ to $-8\%$, declarable on five tasks): the extra energy
is spent on the additional vertical travel of the feet, not on holding larger
joint loads.

\paragraph{The composition is not the sum.} \Cth{} is \Ct{} with the attitude
loop in series after the terrain search, so the $+279\%$ pitch of the search
and the $-82\%$ of the attitude loop are measured on the same signal path.
Comparing the two ends, \Cth{} ends up \emph{better than \Co} on pitch
($0.356$ vs $0.514$ deg on flat), i.e.\ the attitude loop more than pays back
the disturbance the search introduced. Whether a controller consisting of
attitude feedback \emph{without} terrain search would do better still is the
obvious next ablation, and we did not run it: the plant for it does not exist,
and building it would have meant modifying a model that was already frozen for
the campaign.

\subsection{A finding about the composition itself: stance detection degrades}

The clearest architectural result of the project was not planned. When the
ramp task (T4) was analysed on \Cth, the kinematic stance detector
(Section~\ref{sec:setup}) --- a foot counts as in stance when its world speed
stays under \SI{0.05}{\meter\per\second} for at least a quarter of the nominal
stance --- found essentially no stance at all. Our first reading was that the
robot was not walking. It was wrong, and the record of why is worth keeping:
the mean body height we quoted was the ramp's own height, and the
force-sensor stance count of $1.00$ is meaningless on a ramp, because those
sensors only see the floor plane.

The correct reading came from counting, per leg, the fraction of samples under
threshold and the longest consecutive run:

\begin{center}
\small
\begin{tabular}{@{}lrr@{}}
\toprule
 & samples under \SI{0.05}{\meter\per\second} & longest consecutive stretch \\
\midrule
\Ct  & 41--54\% & 71--113 samples ($\approx$ one full stance) \\
\Cth & 4--8\%   & 3--12 samples \\
\bottomrule
\end{tabular}
\end{center}

The feet are not failing to land: they are landing and then being moved again.
The attitude loop adds a continuous height correction
$\Delta z = x\,u_\theta + y\,u_\phi$ to every foot, stance feet included, so a
foot in stance is never at rest in the world frame. The amplitude is small ---
with a median foot speed of \SI{0.14}{\meter\per\second} the motion over one
fragmented stance is about \SI{17}{\milli\meter}, contained by the saturation
and the rate limiter --- so the gait is unaffected. What breaks is the
\emph{observability of contact by any rest-based criterion}.

This matters beyond our test bench, because the terrain search underneath is
itself a contact-detection mechanism. Composing an attitude loop in series
after a contact-triggered mechanism feeds the second one a world in which
nothing is ever stationary. The composition works here because the terrain
search uses a torque residual rather than foot kinematics, and because the
attitude correction is rate-limited; it would not survive a stance detector
built on foot velocity. We consider this the most transferable thing the
project produced.

\subsection{What the task battery could not settle}

Three of the seven tasks returned no verdict, and saying so is part of the
result.

\begin{itemize}
\item \textbf{T6 (seven-obstacle course) does not discriminate.} None of the
three controllers completes it: $2.91$, $3.12$, \SI{2.99}{\meter} travelled,
all three below the validity threshold on lateral drift
(\SI{0.183}{\meter}, derived from the course geometry, not from the
controllers). Beyond that drift the robot may have gone \emph{around}
obstacle~3 rather than over it, so neither the distance nor any downstream
metric can be compared: the three controllers did not meet the same obstacles.
This is a change from an earlier campaign, before the slew-rate limiter was
present on all three plants, in which \Ct{} appeared to reach \SI{3.71}{\meter}
against \SI{2.44}{\meter} for \Co. That headline is superseded and should not
be quoted.

\item \textbf{T7 (lateral impulse) separates only at the small amplitude.} At
$1\times$ the residual deviation is $1.34$, $6.04$, \SI{0.164}{\milli\meter}
--- a $97\%$ improvement for \Cth{} --- but the noise floor was not measured on
this task, so no verdict is declared. At $16\times$ all three end up at
$392$--\SI{400}{\milli\meter} and none recovers: the disturbance is simply
larger than any of the three can reject.

\item \textbf{Peak torque never clears the $3\times$ rule}, on any task or any
comparison, despite differences as large as $+87\%$. Peak is a single-sample
statistic of an impact and its run-to-run scatter is of the same order as the
differences between controllers. The RMS column, which is an integral, is
declarable on five tasks. We report both and draw conclusions only from the
second.
\end{itemize}

% =====================================================================
\section{Limitations}
\label{sec:limits}

\begin{description}[style=unboxed,leftmargin=0pt,font=\normalfont\itshape]

\item[Ideal actuators.] The joints are commanded in position and the simulator
supplies whatever torque the command requires. Section~\ref{sec:feasibility}
therefore measures \emph{demand}, not achievable behaviour. A campaign with
the AX-12A's speed and torque limits in the loop would likely compress the
differences, because \Cth{} is the controller that asks for the most
additional foot travel per unit time.

\item[The contact sensors see only the floor.] The Simscape force sensors are
attached to the ground plane, so on the ramp, the bump and the obstacles they
report nothing. Every stance-dependent metric on T4, T4D, T5 and T6 comes from
the kinematic detector instead --- which, as Section~\ref{sec:discussion}
shows, is itself degraded on \Cth. Where the two detectors disagree we report
the disagreement rather than picking one.

\item[The attitude lever arms are hardcoded.] The $x$ and $y$ used in
$\Delta z = x\,u_\theta + y\,u_\phi$ are constants in the model, about
$1.8\times$ smaller than the nominal foot positions \texttt{cfg.pf\_nom},
though with correct signs and leg ordering. We chose \emph{not} to correct
them: the proportional gains were tuned empirically with those values in the
loop, so changing the lever arms without re-tuning would change the loop gain
by the same factor. The pair (lever arm, gain) is consistent; only its
factorisation is arbitrary.

\item[The inverse-dynamics estimator is not the plant.] After the correction of
Section~\ref{sec:estimator} the residual mismatch is the \SI{1.6}{\percent}
of mass in the feet, which the URDF tree does not carry. The contact threshold
$\sigma_j$ absorbs it, but it was tuned against this bench and would have to be
re-tuned on another.

\item[The threshold is self-tuned.] $\sigma_j$ was set from the residual
observed during known-free swing on this plant. It is therefore not an
independent parameter, and the terrain search cannot be credited with
robustness to a threshold it was fitted to.

\item[Quasi-static stability only.] The Froude number of the nominal march is
$v^2/(gh) = 0.0095$, so the ZMP collapses onto the ground projection of the
centre of mass to within the precision of the measurement; the stability
analysis assumes a planar contact patch and reports the stance-plane tilt so
that the assumption can be checked rather than trusted.

\item[No hardware.] Nothing in this report was measured on a physical
PhantomX. The platform, the URDF and the servo datasheet are real; the
comparison is not validated.

\end{description}

% =====================================================================
\section{Contributions by team member}
\label{sec:contributions}

\noindent\fcolorbox{red!50}{red!3}{\begin{minipage}{0.97\linewidth}
\small\textbf{TODO --- to be confirmed and completed by the authors before
submission.} The split below reflects the division of work as it actually
happened; names and any item either of us wants to reclaim must be fixed by
hand.
\end{minipage}}

\vspace{6pt}

\begin{description}[style=unboxed,leftmargin=0pt,font=\normalfont\bfseries]

\item[Student 1 --- plant and controllers.]
Construction of the Simscape Multibody plant from the PhantomX URDF;
implementation of the open-loop tripod generator (\Co); implementation of the
terrain-search layer after Arrigoni et al.\ (\Ct), including the
inverse-dynamics residual and the per-leg contact flag; design,
implementation and empirical tuning of the attitude loop (\Cth) --- pitch
first, then roll --- and of the saturation and rate limiter that bound its
authority. Ownership of the Simulink models throughout.

\item[Student 2 --- measurement, diagnosis and analysis.]
Design of the seven-task battery and of the measurement protocol
(Section~\ref{sec:method}), including the noise-floor procedure and the
$3\times$ declarability rule; the full measurement campaign and its automation
(single-point controller selection, batch execution, result conversion,
comparison table generation); identification, quantification and correction of
the two test-bench defects of Section~\ref{sec:defects}; the stance-detection
diagnosis of Section~\ref{sec:discussion}; the quasi-static stability check;
the feasibility analysis against the servo datasheet; repository structure and
documentation.

\item[Joint.]
Choice of the three controllers and of the ablation structure; reconstruction
and justification of the gait parameters (Section~\ref{sec:setup});
interpretation of the results; this report.

\end{description}

% =====================================================================
\section{Relation to other coursework and novelty}
\label{sec:novelty}

This project was not derived from, and does not extend, a project submitted
for another course. The platform model, the three controllers, the task
battery, the measurement protocol and all the results reported here were
produced for Field and Service Robotics.

Relative to the course material, the project is an application rather than an
extension of the theory, and its novelty lies in three places:

\begin{enumerate}
\item \textbf{A controlled ablation instead of a demonstration.} The usual
shape of a project of this kind is: implement a controller, show that it
walks. Here the same plant, ground, contact law and gait generator carry three
controllers that differ by exactly one mechanism each, so each comparison
isolates one mechanism. That design is what makes it possible to say that
terrain search costs attitude and attitude feedback costs energy, rather than
that one full controller is better than another.

\item \textbf{A declarability criterion fixed before the campaign.} Simulation
results are reproducible to the last digit if nothing is perturbed, which
makes every difference look significant. Perturbing the initial height by
\SI{\pm 0.2}{\milli\meter} over three runs gives a noise floor, and requiring a
difference to exceed $3\times$ the worse of the two floors before it is
declared converts a table of numbers into a table of claims. Roughly a third
of the differences we measured do not survive that rule, including every
peak-torque comparison.

\item \textbf{Two silent bench defects, found by measurement rather than by
inspection.} Both (Section~\ref{sec:defects}) produced plausible,
publishable-looking numbers while being wrong, and neither was visible in the
block diagram. The estimator-misalignment one in particular is a direct
instance of the course's point about model-based residuals: correcting the
plant without correcting the model that predicts it is worse than correcting
neither, because the residual saturates and the contact flag degenerates to a
constant.
\end{enumerate}

% =====================================================================
\section{Conclusions}
\label{sec:conclusions}

Three controllers for a PhantomX hexapod were compared on one plant over seven
tasks, under a protocol that declares a difference only when it exceeds three
times the measured noise floor.

\begin{itemize}
\item The terrain search of~\cite{arrigoni2024}, taken on its own, is a net
cost on terrain the gait generator already models: pitch worsens by
$219$--$279\%$ with no energy return. On irregular terrain it improves roll
($-30$ to $-49\%$), cost of transport ($-7$ to $-10\%$) and speed
($+2$ to $+3\%$), while still worsening pitch.
\item The attitude loop is the mechanism that buys attitude, and it does so on
every task: pitch $-73$ to $-82\%$ at $18$ to $94$ times the noise floor, and
on the $8^\circ$ ramp a qualitatively different behaviour (trunk held
horizontal rather than parallel to the slope). It costs $28$--$40\%$ in energy
and cost of transport, spent on foot travel rather than on joint load --- the
continuous torque actually falls.
\item Composed in series, the two mechanisms leave the body better than the
open-loop baseline on pitch, but they make foot contact unobservable to any
rest-based detector: stance stretches collapse from $71$--$113$ samples to
$3$--$12$.
\item On the real servos, the march itself would run at $39$--$46\%$ of stall
on the worst-loaded joint --- $2.3\times$ the manufacturer's recommended
continuous load --- with impact peaks above stall on at most $0.91\%$ of
samples.
\end{itemize}

The obvious continuation is the fourth cell of the ablation --- attitude
feedback without terrain search --- which would say whether the search
contributes anything at all once the body is stabilised, and a campaign with
the servo limits in the loop, which is where the energy cost of \Cth{} would
stop being free.

% =====================================================================
\begin{thebibliography}{9}

\bibitem{arrigoni2024}
S.~Arrigoni, F.~Braghin, et al.,
``Control of a Hexapod Robot Considering Terrain Interaction,''
\emph{Robotics}, vol.~13, no.~10, art.~142, 2024.
% TODO: verificare la lista completa degli autori sul PDF

\bibitem{ax12a}
ROBOTIS, \emph{Dynamixel AX-12A e-Manual}, ROBOTIS Co., Ltd.
\url{https://emanual.robotis.com/docs/en/dxl/ax/ax-12a/}

\bibitem{phantomx}
Trossen Robotics, \emph{PhantomX AX Metal Hexapod Mark~II --- assembly and
specifications}.

\bibitem{simscape}
The MathWorks, Inc., \emph{Simscape Multibody User's Guide}, R2025b.

\end{thebibliography}

\end{document}
